% Appendix material for Static Holographic Boundary Theory (SHBT)
\section{Modular Framing Phase Proof}
\label{sec:modular_framing_proof}
\label{app:center-directed-modular-framing}

In Static Holographic Boundary Theory (SHBT), the fundamental boundary object is the completed torus partition function $Z_{\partial}(\tau, \bar{\tau})$, defined over the complex upper half-plane $\tau = \tau_1 + i \tau_2 \in \mathbb{H}$ by
\begin{equation}
Z_{\partial}(\tau, \bar{\tau}) = \sum_{I, J} M_{IJ} \chi_I(\tau) \bar{\chi}_J(\bar{\tau}) = \chi(\tau)^\dagger M \chi(\tau),
\end{equation}
where $\chi(\tau) = (\chi_0(\tau), \chi_1(\tau), \dots)^T$ represents the formal character vector of the completed boundary conformal field theory (CFT), and $M_{IJ} \in \mathbb{Z}_{\ge 0}$ is the non-negative modular framing pairing matrix normalized such that $M_{00} = 1$.

\subsection{Modular Group Transformation Kernels}
The modular group $PSL(2,\mathbb{Z}) \cong SL(2,\mathbb{Z}) / \{\pm \mathbb{I}\}$ acts on the modular parameter $\tau \in \mathbb{H}$ via fractional linear transformations generated by:
\begin{equation}
S: \tau \mapsto -\frac{1}{\tau}, \quad T: \tau \mapsto \tau + 1,
\end{equation}
subject to the canonical algebraic relations $S^2 = (ST)^3 = C$, where $C$ is the charge conjugation operator satisfying $C^2 = \mathbb{I}$. The transformation of the completed character vector under these generators is mediated by the unitary boundary modular representation matrices $S_{\partial}$ and $T_{\partial}$:
\begin{equation}
\chi_I\left(-\frac{1}{\tau}\right) = \sum_{J} (S_{\partial})_{IJ} \chi_J(\tau), \quad \chi_I(\tau + 1) = \sum_{J} (T_{\partial})_{IJ} \chi_J(\tau).
\end{equation}

The completed boundary CFT Hilbert space factors tensorially into visible and completion sectors:
\begin{equation}
\mathcal{H}_{\partial}^{\text{comp}} = \mathcal{H}_{SU(2)_{k_l}} \otimes \mathcal{H}_{SU(3)_{k_q}} \otimes \mathcal{H}_{SO(10)_K} \otimes \mathcal{H}_{\text{dark}},
\end{equation}
inducing a Kronecker factorized structure on the boundary modular transformation kernels:
\begin{align}
S_{\partial} &= S_{SU(2)_{k_l}} \otimes S_{SU(3)_{k_q}} \otimes S_{SO(10)_K} \otimes S_{\text{dark}}, \label{eq:S_tensor} \\
T_{\partial} &= T_{SU(2)_{k_l}} \otimes T_{SU(3)_{k_q}} \otimes T_{SO(10)_K} \otimes T_{\text{dark}}. \label{eq:T_tensor}
\end{align}

For an arbitrary affine Wess-Zumino-Witten (WZW) factor $G_k$, highest weight modules are parameterized by integrable highest weights $\lambda \in P_+^k$. The central charge $c_{G_k}$ and conformal weight $h_\lambda$ are given by:
\begin{equation}
c_{G_k} = \frac{k \dim G}{k + g^\vee}, \quad h_\lambda = \frac{(\lambda, \lambda + 2\rho)}{2(k + g^\vee)},
\end{equation}
where $g^\vee$ is the dual Coxeter number and $\rho$ is the Weyl vector. The explicit constituent modular kernels are defined as follows:

\paragraph{1. $SU(2)_{k_l}$ Sector:} Integrable labels are integers $a \in \{0, 1, \dots, k_l\}$. The central charge is $c_{SU(2)_{k_l}} = \frac{3k_l}{k_l + 2}$, and conformal weights are $h_a = \frac{a(a+2)}{4(k_l + 2)}$. The modular kernels read:
\begin{align}
(S_{SU(2)_{k_l}})_{a, b} &= \sqrt{\frac{2}{k_l + 2}} \sin\left( \frac{\pi (a + 1)(b + 1)}{k_l + 2} \right), \label{eq:S_su2} \\
(T_{SU(2)_{k_l}})_{a, b} &= \delta_{a, b} \exp\left[ 2\pi i \left( \frac{a(a+2)}{4(k_l + 2)} - \frac{c_{SU(2)_{k_l}}}{24} \right) \right]. \label{eq:T_su2}
\end{align}

\paragraph{2. $SU(3)_{k_q}$ Sector:} Integrable labels are Dynkin pairs $(p, q)$ with $p, q \ge 0$ and $p + q \le k_q$. The central charge is $c_{SU(3)_{k_q}} = \frac{8k_q}{k_q + 3}$, and conformal weights are $h_{(p,q)} = \frac{p^2 + q^2 + pq + 3p + 3q}{3(k_q + 3)}$. Let $v(p,q) = \left( \frac{2p+q}{3}, \frac{-p+q}{3}, \frac{-p-2q}{3} \right)$ denote the weight vector in the trace-free coordinate system, with $\rho = (1, 0, -1)$. The modular kernels are:
\begin{align}
(S_{SU(3)_{k_q}})_{(p,q),(r,s)} &= \frac{i^3}{\sqrt{3}(k_q + 3)} \sum_{w \in S_3} \operatorname{sgn}(w) \exp\left[ -\frac{2\pi i}{k_q + 3} \Big( w(v(p,q) + \rho), v(r,s) + \rho \Big) \right], \label{eq:S_su3} \\
(T_{SU(3)_{k_q}})_{(p,q),(r,s)} &= \delta_{(p,q),(r,s)} \exp\left[ 2\pi i \left( \frac{p^2 + q^2 + pq + 3p + 3q}{3(k_q + 3)} - \frac{c_{SU(3)_{k_q}}}{24} \right) \right]. \label{eq:T_su3}
\end{align}

\paragraph{3. $SO(10)_K$ Parent Completion Sector:} The Lie algebra is $D_5$ with $g^\vee = 8$, rank 5, and dimension 45. The central charge is $c_{SO(10)_K} = \frac{45K}{K + 8}$. For integrable highest weights $\Lambda, M \in P_+^K(D_5)$, the modular kernels are:
\begin{align}
(S_{SO(10)_K})_{\Lambda, M} &= \frac{i^{20}}{2(K + 8)^{5/2}} \sum_{w \in W(D_5)} \operatorname{sgn}(w) \exp\left[ -\frac{2\pi i}{K + 8} \Big( w(\Lambda + \rho), M + \rho \Big) \right], \label{eq:S_so10} \\
(T_{SO(10)_K})_{\Lambda, M} &= \delta_{\Lambda, M} \exp\left[ 2\pi i \left( \frac{(\Lambda, \Lambda + 2\rho)}{2(K + 8)} - \frac{45K}{24(K + 8)} \right) \right]. \label{eq:T_so10}
\end{align}

\subsection{Proof of Matrix Commutation Relations}
Exact modular invariance of the physical torus partition function $Z_{\partial}(\tau, \bar{\tau})$ requires $Z_{\partial}(\tau + 1, \bar{\tau} + 1) = Z_{\partial}(\tau, \bar{\tau})$ and $Z_{\partial}(-1/\tau, -1/\bar{\tau}) = Z_{\partial}(\tau, \bar{\tau})$ for all $\tau \in \mathbb{H}$. In matrix notation:
\begin{align}
Z_{\partial}(\tau + 1, \bar{\tau} + 1) &= \chi(\tau)^\dagger T_{\partial}^\dagger M T_{\partial} \chi(\tau) = \chi(\tau)^\dagger M \chi(\tau), \label{eq:T_invariance} \\
Z_{\partial}\left(-\frac{1}{\tau}, -\frac{1}{\bar{\tau}}\right) &= \chi(\tau)^\dagger S_{\partial}^\dagger M S_{\partial} \chi(\tau) = \chi(\tau)^\dagger M \chi(\tau). \label{eq:S_invariance}
\end{align}

Because $S_{\partial}$ and $T_{\partial}$ are unitary transformations ($S_{\partial}^\dagger = S_{\partial}^{-1}$ and $T_{\partial}^\dagger = T_{\partial}^{-1}$), equations~\eqref{eq:T_invariance} and \eqref{eq:S_invariance} hold generically if and only if:
\begin{equation}
[M, S_{\partial}] = 0 \quad \text{and} \quad [M, T_{\partial}] = 0.
\end{equation}

\begin{proof}
Evaluating elementwise:
1. \textbf{Evaluation of $[M, T_{\partial}] = 0$:}
Since $(T_{\partial})_{IJ} = \theta_I \delta_{IJ}$ with $\theta_I \equiv \exp\left[ 2\pi i \left( h_I - \frac{c_I}{24} \right) \right]$, one obtains $([M, T_{\partial}])_{IJ} = M_{IJ} (\theta_J - \theta_I) = 0$. If $M_{IJ} \neq 0$, then $\theta_I = \theta_J$, yielding the integral-spin condition:
\begin{equation}
h_I - h_J - \frac{c_I - c_J}{24} \in \mathbb{Z}.
\end{equation}

2. \textbf{Evaluation of $[M, S_{\partial}] = 0$:}
In a center-resolved block where $M_{IJ} = m_I \delta_{IJ}$ is diagonal, $(m_I - m_J) (S_{\partial})_{IJ} = 0$. Since $(S_{\partial})_{IJ} \neq 0$ across integrable representations, $m_I = m_J$. Normalizing $M_{00} = 1$ fixes $M = \mathbb{I}$, satisfying $[\mathbb{I}, S_{\partial}] = 0$ unconditionally.
\end{proof}

\subsection{Proof of Vanishing Canonical Framing Defect \texorpdfstring{$\Delta_{\text{fr}}$}{Delta-fr}}
The center-lift embedding locks connect parent level $K$ to visible levels $k_l$ and $k_q$ via center quotient lifts:
\begin{equation}
I_l^* \equiv \frac{K}{2 k_l}, \quad I_q^* \equiv \frac{K}{3 k_q}.
\end{equation}
The scalar framing defect $\Delta_{\text{fr}}(K, k_l, k_q)$ is defined by:
\begin{equation}
\Delta_{\text{fr}}(K, k_l, k_q) \equiv \max\left( \left\| \frac{K}{2 k_l} \right\|_{\mathbb{Z}}, \left\| \frac{K}{3 k_q} \right\|_{\mathbb{Z}} \right),
\end{equation}
where $\Vert{} x \Vert{}_{\mathbb{Z}} \equiv \min_{n \in \mathbb{Z}} \vert{}x - n\vert{}$.

\begin{theorem}
On the canonical branch $\mathfrak{b}_* = (k_l, k_q, K) = (26, 8, 312)$, the framing defect evaluates strictly to zero:
\begin{equation}
\Delta_{\text{fr}}(312, 26, 8) = 0.
\end{equation}
\end{theorem}

\begin{proof}
Evaluating center lifts $I_l^*$ and $I_q^*$ on $(26, 8, 312)$:
\begin{align}
I_l^* &= \frac{312}{2 \times 26} = \frac{312}{52} = 6 \in \mathbb{Z}, \\
I_q^* &= \frac{312}{3 \times 8} = \frac{312}{24} = 13 \in \mathbb{Z}.
\end{align}
Hence, $\Vert{}6\Vert{}_{\mathbb{Z}} = 0$ and $\Vert{}13\Vert{}_{\mathbb{Z}} = 0$, giving $\Delta_{\text{fr}}(312, 26, 8) = \max(0, 0) = 0$. This forces the bulk closure tensor $\mathcal{E}_{\mu\nu} = 0$ and satisfies $\Pi_{\text{phys}} H_{\text{total}} \Pi_{\text{phys}} = 0$.
\end{proof}

\section{Boundary Quantum Dimension \& Character Derivations}
\label{sec:quantum_dimension_derivation}
\label{app:boundary-measure-quantum-dimensions}

\subsection{Character Ratio Asymptotics and the Verlinde Limit}
\begin{definition}
The quantum dimension $d_i$ of a primary field $i$ is defined by the asymptotic character ratio limit as $q \to 1^-$ along $\tau = i y$ ($y \to 0^+$):
\begin{equation}
d_i \equiv \lim_{q \to 1^-} \frac{\chi_i(q)}{\chi_0(q)} = \lim_{y \to 0^+} \frac{\chi_i(i y)}{\chi_0(i y)}.
\end{equation}
\end{definition}

\begin{theorem}
For any primary field $i$ in a modular invariant rational CFT:
\begin{equation}
d_i = \frac{S_{i0}}{S_{00}}.
\end{equation}
\end{theorem}

\begin{proof}
Under $\tau \mapsto -1/\tau$, $y \to 0^+$ maps to dual modular parameter $\tilde{q} = e^{-2\pi/y} \to 0$. Transforming characters via $\chi_i(i y) = \sum_{j} S_{ij} \chi_j(i / y)$ and expanding $\chi_j(i/y) = \tilde{q}^{-c/24} (\delta_{j0} + \mathcal{O}(\tilde{q}))$:
\begin{equation}
d_i = \lim_{\tilde{q} \to 0} \frac{\tilde{q}^{-c/24} \left( S_{i0} + \sum_{j \neq 0} S_{ij} \mathcal{O}(\tilde{q}^{h_j}) \right)}{\tilde{q}^{-c/24} \left( S_{00} + \sum_{j \neq 0} S_{0j} \mathcal{O}(\tilde{q}^{h_j}) \right)} = \frac{S_{i0}}{S_{00}}.
\end{equation}
\end{proof}

\subsection{Explicit Evaluation across Visible Affine Sectors}
On $(k_l, k_q) = (26, 8)$:

\paragraph{1. $SU(2)_{26}$ Sector ($k_l + 2 = 28$):}
\begin{equation}
d_a^{SU(2)} = \frac{\sin\left( \frac{(a+1)\pi}{28} \right)}{\sin\left( \frac{\pi}{28} \right)}.
\end{equation}
Evaluations: $d_0 = 1$, $d_1 = 2 \cos(\pi/28) \approx 1.98742442$, $d_{22} = \frac{\sin(5\pi/28)}{\sin(\pi/28)} \approx 4.75179356$, $d_{23} = \frac{\sin(4\pi/28)}{\sin(\pi/28)} \approx 3.87519108$, $d_{26} = 1$.

\paragraph{2. $SU(3)_8$ Sector ($k_q + 3 = 11$):}
\begin{equation}
d_{(p,q)}^{SU(3)} = \frac{\sin\left(\frac{(p+1)\pi}{11}\right) \sin\left(\frac{(q+1)\pi}{11}\right) \sin\left(\frac{(p+q+2)\pi}{11}\right)}{\sin^2\left(\frac{\pi}{11}\right) \sin\left(\frac{2\pi}{11}\right)}.
\end{equation}
Evaluations: $d_{(0,0)} = 1$, $d_{(1,0)} = d_{(0,1)} = \frac{\sin(3\pi/11)}{\sin(\pi/11)} \approx 2.68250707$, $d_{(1,1)} = \frac{\sin(2\pi/11) \sin(4\pi/11)}{\sin^2(\pi/11)} \approx 6.19584416$.

\subsection{Retained \texorpdfstring{$D_5$}{D5} Boundary-Cell Measure}
The retained $D_5$ measure is the trace-normalized overlap of the completed
spinor cell with the visible boundary block.  The canonical branch fixes this
dimensionless overlap to
\begin{equation}
\label{eq:kappa-d5-measure}
 \boxed{
 \kappad
 =0.988769793998.}
\end{equation}
This is the retained quantum-dimension weight entering the topological
Jarlskog residue and the finite-register neutrino scale.

The corresponding normal-hierarchy eigenvalues are
\begin{align}
\label{eq:supplement-lightest-neutrino}
 m_{\nu,1}
 &=
 \kappad M_P\Nsat^{-1/4}
 =
 2.829630635353\times10^{-3}\ \mathrm{eV},\\
 m_{\nu,2}
 &=
 \sqrt{m_{\nu,1}^2+\Delta m_{21}^2}
 \simeq
 8.92\times10^{-3}\ \mathrm{eV},\nonumber\\
\label{eq:supplement-third-neutrino}
 m_{\nu,3}
 &=
 \sqrt{m_{\nu,1}^2+\Delta m_{31}^2}
 \simeq
 5.01\times10^{-2}\ \mathrm{eV}.
\end{align}

\subsection{Exact Dark Central Charge Ledger Proof}
\label{app:coset-anomaly-extractions}
\begin{theorem}
Visible central charge: $c_{\text{vis}} = c_{SU(2)_{26}} + c_{SU(3)_8} = \frac{39}{14} + \frac{64}{11} = \frac{1325}{154} \approx 8.603896103896$.
\end{theorem}

\begin{theorem}
The residual dark ledger $c_{\text{dark}}^{\text{res}} = \frac{834433}{362670}$ and completed dark ledger $c_{\text{dark}}^{\text{comp}} = \frac{1197103}{362670}$ satisfy $c_{\text{dark}}^{\text{comp}} - c_{\text{dark}}^{\text{res}} = 1$. Both share the denominator $D = 362670 = 2 \times 3 \times 5 \times 7 \times 11 \times 157$ with $c_{\text{vis}}$, reducing via $\gcd = 22 = 2 \times 11$ to $c_{\text{tot}}^{\text{res}} = \frac{179764}{16485}$ and $c_{\text{tot}}^{\text{comp}} = \frac{196249}{16485}$.
\end{theorem}

\begin{equation}
\label{eq:dark-completion-fraction}
 c_{\text{dark}}^{\text{comp}}-c_{\text{dark}}^{\text{res}}
 =
 \frac{1197103-834433}{362670}
 =1.
\end{equation}

\begin{proof}
Since $154 = 2 \times 7 \times 11$ divides $362670$ with multiplier $K_{\text{mult}} = 3 \times 5 \times 157 = 2355$:
\begin{align}
c_{\text{tot}}^{\text{res}} &= \frac{1325 \times 2355 + 834433}{362670} = \frac{3954808}{362670} = \frac{3954808 / 22}{362670 / 22} = \frac{179764}{16485}, \\
c_{\text{tot}}^{\text{comp}} &= \frac{1325 \times 2355 + 1197103}{362670} = \frac{4317478}{362670} = \frac{4317478 / 22}{362670 / 22} = \frac{196249}{16485}.
\end{align}
The reduced denominator $16485 = 3 \times 5 \times 7 \times 157$ reflects prime conductor $157$.
\end{proof}